\section{The Five Parameters and Their Ranges}
\label{sec:parameters}

The \bisect{} algorithm exposes five parameters that a deploying authority
must choose.
Table~\ref{tab:params} lists the parameters, their technical roles, default
values, and admissible ranges.

\begin{table}[h]
\centering
\caption{The five \bisect{} parameters. Default is the DIA-recommended value.
         Admissible range is the technically valid range for which the algorithm
         produces legally compliant districts (contiguous, population-balanced,
         VRA-compliant).}
\label{tab:params}
\begin{tabular}{llccl}
\toprule
Parameter & Symbol & Default & Admissible range & Technical role \\
\midrule
METIS imbalance & $\ufactor$ & 1.0\% & 0.1\%--5.0\% & Population balance tolerance \\
County edge weight & $\acounty$ & 3.0 & 1.0--10.0 & County integrity preservation \\
Convergence threshold & $T$ & 600 & 100--1{,}000 & \cs{} stopping criterion \\
Area swing & $\aswing$ & 1.10 & 1.05--1.20 & Area/compactness tradeoff \\
VRA boost weight & $\wvra$ & 0.40 & 0.2--0.6 & Minority community protection \\
\bottomrule
\end{tabular}
\end{table}

\subsection{Parameter Interpretations}

\paragraph{$\ufactor$ (METIS imbalance factor).}
METIS's $k$-way partition algorithm accepts a population imbalance tolerance.
At $\ufactor = 1\%$, no district may deviate more than 1\% from ideal
population.
At $\ufactor = 5\%$, deviations up to 5\% are permitted.
Higher $\ufactor$ gives the partitioner more freedom to optimise compactness
at the cost of population equality.
Federal courts require districts to be ``as equal as practicable'';
$\ufactor \leq 1\%$ is the standard for congressional plans.

\paragraph{$\acounty$ (county edge weight).}
Edges between census tracts in the same county receive weight $\acounty$
(default: 3) relative to cross-county edges (weight: 1).
Higher $\acounty$ encourages the bisection to avoid splitting counties.
At $\acounty = 1$, county membership is ignored; at $\acounty = 10$,
county boundaries are strongly preferred.

\paragraph{$T$ (convergence sweep threshold).}
The \cs{} algorithm~\citep{b16paper} sweeps seeds until $T$ consecutive
seeds produce no improvement in the objective function.
Higher $T$ gives a more thorough search at the cost of longer runtime.
At $T = 100$, the algorithm occasionally misses the global optimum
(B.7 shows the tail extends to 511 consecutive non-improving seeds for Georgia).
At $T = 1{,}000$, runtime is nearly $10\times$ the default.

\paragraph{$\aswing$ (area swing).}
In the GeoSection (ratio-optimal) algorithm, the area tolerance parameter
controls the permitted deviation from the ideal area-to-population ratio
when splitting a region.
Higher $\aswing$ permits more area imbalance in exchange for better
population balance.

\paragraph{$\wvra$ (VRA boost weight).}
Edges adjacent to VRA-eligible minority communities receive an additional
weight $\wvra$ to encourage the bisection to keep minority communities
intact.
Higher $\wvra$ provides stronger VRA protection at some compactness cost.
